\documentclass[11pt,a4paper]{article}

\usepackage[utf8]{inputenc}
\usepackage[T1]{fontenc}
\usepackage[french]{babel}
\usepackage{amsmath}
\usepackage{amssymb}
\usepackage{booktabs}
\usepackage{geometry}
\usepackage{hyperref}
\usepackage{longtable}
\usepackage{enumitem}

\geometry{margin=2.3cm}
\hypersetup{
  colorlinks=true,
  linkcolor=blue,
  urlcolor=blue
}

\title{Detail mathematique du pipeline d'analyse DXD}
\author{}
\date{}

\newcommand{\RMS}{\operatorname{RMS}}
\newcommand{\median}{\operatorname{median}}
\newcommand{\IQR}{\operatorname{IQR}}
\newcommand{\MAD}{\operatorname{MAD}}
\newcommand{\finite}{\mathcal{F}}

\begin{document}
\maketitle

\section{Objet du document}

Ce document decrit les calculs effectues par les scripts du depot, de la lecture des fichiers DXD jusqu'aux index exploites par l'application.
Il suit principalement le flux moderne :

\begin{enumerate}[label=\arabic*.]
  \item \texttt{00\_scan\_dxd\_channel\_presence.py} : lecture DXD, resampling, export JSON canonique, features globales et glissantes.
  \item \texttt{build\_window\_quality\_index.py} : calcul par fenetre des residus physiques, scores qualite et labels de contexte.
  \item \texttt{build\_window\_cluster\_index.py} : regroupement DBSCAN de fenetres dynamiquement comparables.
  \item \texttt{build\_speed\_day\_summary.py}, \texttt{build\_drift\_index.py}, \texttt{build\_correlation\_anomaly\_index.py} : agregations et index secondaires.
\end{enumerate}

Les scripts \texttt{01\_export\_dxd\_resampled\_json.py} et \texttt{02\_extract\_features.py} constituent un flux historique ou auxiliaire. Leurs calculs sont resumes en fin de document.

\section{Notations generales}

Pour une acquisition, les signaux sont resamples sur une grille temporelle commune :
\[
  t_i = t_{\min} + i \Delta t,\qquad i=0,\ldots,N-1.
\]
Dans la configuration par defaut de l'export DXD :
\[
  f_s = 100~\mathrm{Hz}, \qquad \Delta t = \frac{1}{f_s}=0.01~\mathrm{s}.
\]

Les grandeurs principales sont :
\[
\begin{aligned}
  v_x(t) &: \text{vitesse longitudinale vehicule en m/s},\\
  v_y(t) &: \text{vitesse laterale vehicule en m/s},\\
  a_x(t) &: \text{acceleration longitudinale en } \mathrm{m/s^2},\\
  a_y(t) &: \text{acceleration laterale en } \mathrm{m/s^2},\\
  r(t) &: \text{vitesse de lacet, ou yaw rate, en rad/s},\\
  \omega_{fl},\omega_{fr},\omega_{rl},\omega_{rr} &: \text{vitesses angulaires des roues}.
\end{aligned}
\]

Les valeurs non finies (\texttt{NaN}, infinis) sont ignorees dans les moyennes, RMS et quantiles lorsque le code utilise les fonctions \texttt{safe\_*}. On note :
\[
  \finite(x)=\{i \mid x_i \text{ est fini}\}.
\]

Pour un vecteur $x$, les statistiques de base utilisees sont :
\[
  \overline{x}=\frac{1}{|\finite(x)|}\sum_{i\in\finite(x)}x_i,
  \qquad
  \RMS(x)=\sqrt{\frac{1}{|\finite(x)|}\sum_{i\in\finite(x)}x_i^2},
\]
\[
  \max |x| = \max_{i\in\finite(x)} |x_i|,
  \qquad
  q_{95}(|x|)=\text{quantile a 95\% de } |x_i|.
\]

La derivee numerique est calculee par gradient sur les points finis :
\[
  \dot{x}_i \simeq \nabla x_i / \Delta t.
\]

\section{Normalisation des canaux}

Le module \texttt{dxd\_schema.py} transforme les noms sources Dewesoft en noms canoniques. Exemples :

\begin{longtable}{lll}
\toprule
Nom source & Nom canonique & Conversion \\
\midrule
\texttt{VelX (km/h)} & \texttt{vehicle.vx} & $v_x = \text{VelX}/3.6$ \\
\texttt{VelY (km/h)} & \texttt{vehicle.vy} & $v_y = \text{VelY}/3.6$ \\
\texttt{Vel (km/h)} & \texttt{vehicle.speed} & $v = \text{Vel}/3.6$ \\
\texttt{VehYaw\_W\_Actl (rad/s)} & \texttt{vehicle.yaw\_rate} & identite \\
\texttt{VehLong\_A\_Actl (m/s\^{}2)} & \texttt{vehicle.ax} & identite \\
\texttt{VehLat\_A\_Actl (m/s\^{}2)} & \texttt{vehicle.ay} & identite \\
\texttt{Latitude (\_)} & \texttt{gps.latitude} & identite \\
\texttt{Longitude (\_)} & \texttt{gps.longitude} & identite \\
\bottomrule
\end{longtable}

En general, si un canal source $y^{raw}$ a une echelle $s$ et un offset $b$, le canal canonique vaut :
\[
  y = s\,y^{raw}+b.
\]

\section{Lecture DXD, nettoyage et resampling}

Pour chaque canal lu dans un fichier DXD, le code :

\begin{enumerate}[label=\arabic*.]
  \item extrait les couples $(t_i,y_i)$ ;
  \item supprime les couples non finis ;
  \item trie par temps croissant ;
  \item supprime les doublons temporels en gardant la premiere occurrence ;
  \item calcule l'intersection temporelle commune a tous les canaux selectionnes :
  \[
    t_{\min}=\max_c t^{(c)}_0,\qquad
    t_{\max}=\min_c t^{(c)}_{\mathrm{fin}}.
  \]
\end{enumerate}

La grille commune est :
\[
  t_k = t_{\min}+k\Delta t,\qquad t_k \le t_{\max}+0.5\Delta t.
\]

Chaque canal est interpole lineairement :
\[
  y_c(t_k)=\operatorname{interp}(t_k; t^{(c)}, y^{(c)}),
\]
avec \texttt{NaN} a gauche et a droite hors du support du canal.

\section{Canaux derives dans l'export JSON}

Le script \texttt{00\_scan\_dxd\_channel\_presence.py} ajoute certains residus derives au JSON.
Ces residus sont calcules sur les noms sources avant normalisation complete.

\subsection{Coherence laterale simple}

Si $a_y$, $v_x$ et $r$ sont disponibles :
\[
  e_{lat}(t)=a_y(t)-v_x(t)r(t).
\]
Le JSON contient aussi un RMS glissant sur une fenetre de duree $T_r=2~\mathrm{s}$ par defaut :
\[
  S_{lat}(i)=
  \sqrt{
    \frac{1}{|\mathcal{W}_i|}
    \sum_{j\in\mathcal{W}_i} e_{lat}(j)^2
  },
\]
ou $\mathcal{W}_i$ est la fenetre retrospective de longueur $T_r$ finissant en $i$.

\subsection{Differences de vitesses de roues}

Les residus de roues exportes dans le JSON sont :
\[
\begin{aligned}
  e_{front}(t) &= \omega_{fl}(t)-\omega_{fr}(t),\\
  e_{rear}(t)  &= \omega_{rl}(t)-\omega_{rr}(t),\\
  e_{left}(t)  &= \omega_{fl}(t)-\omega_{rl}(t),\\
  e_{right}(t) &= \omega_{fr}(t)-\omega_{rr}(t).
\end{aligned}
\]
Chaque residu possede egalement un RMS glissant sur $T_r$.

\subsection{Steer, camber et acceleration longitudinale}

Pour les deux voies de direction et de carrossage :
\[
  e_{\delta}(t)=\delta_{S1}(t)+\delta_{S2}(t),
  \qquad
  e_{\gamma}(t)=\gamma_{S1}(t)-\gamma_{S2}(t).
\]

Pour l'acceleration longitudinale :
\[
  e_{ax}(t)=a_x(t)-\frac{d v_x}{dt}(t).
\]

\section{Features globales et glissantes dans le JSON}

Pour chaque canal brut ou derive $y(t)$, le JSON contient des features globales :
\[
  n_{valid},\quad \overline{y},\quad \sigma_y,\quad \min y,\quad \max y,\quad \RMS(y),\quad E_y.
\]
L'energie est :
\[
  E_y=\sum_{i\in\finite(y)} y_i^2\,\Delta t.
\]

Les features glissantes sont calculees par fenetres de duree $T_w=4~\mathrm{s}$ et pas $T_s=1~\mathrm{s}$ par defaut :
\[
  \mathcal{W}_m=\{i_m,\ldots,i_m+n_w-1\},\qquad
  n_w=\operatorname{round}(T_w/\Delta t).
\]
Pour chaque fenetre, le code stocke :
\[
  \overline{y}_{\mathcal{W}},\quad
  \sigma_{y,\mathcal{W}},\quad
  \RMS(y_{\mathcal{W}}),\quad
  \min_{\mathcal{W}}y,\quad
  \max_{\mathcal{W}}y.
\]

Quelques correlations globales sont aussi calculees sur des paires fixes :
\[
  \rho(x,y)=
  \frac{\sum_i (x_i-\overline{x})(y_i-\overline{y})}
       {\sqrt{\sum_i(x_i-\overline{x})^2}\sqrt{\sum_i(y_i-\overline{y})^2}}.
\]

\section{Fenetrage pour l'index qualite}

Le script \texttt{build\_window\_quality\_index.py} lit les JSON resamples et reconstruit des fenetres de duree $T_w=4~\mathrm{s}$ avec un pas $T_s=1~\mathrm{s}$ :
\[
  n_w=\max(2,\operatorname{round}(T_w/\Delta t)),\qquad
  n_s=\max(1,\operatorname{round}(T_s/\Delta t)).
\]
Les fenetres sont :
\[
  \mathcal{W}_m = \{m n_s,\ldots,m n_s+n_w-1\}.
\]

Pour chaque fenetre, les statistiques dynamiques principales sont :
\[
  \overline{v_x},\quad
  \overline{v_x}^{km/h}=3.6\,\overline{v_x},\quad
  \RMS(a_x),\quad
  \RMS(a_y),\quad
  \RMS(r),\quad
  \RMS(\delta).
\]

Deux ratios de vitesse sont calcules :
\[
  p_{zero}=\frac{1}{|\finite(v_x)|}\sum_{i\in\finite(v_x)}\mathbf{1}(|v_x(i)|<1/3.6),
\]
\[
  p_{low}=\frac{1}{|\finite(v_x)|}\sum_{i\in\finite(v_x)}\mathbf{1}(|v_x(i)|<5/3.6).
\]

\section{Qualite de signal par canal}

Pour chaque canal de qualite, dans une fenetre $\mathcal{W}$ :
\[
  r_{finite}=\frac{\#\{i\in\mathcal{W}\mid y_i\text{ fini}\}}{|\mathcal{W}|}.
\]
Le score de presence vaut :
\[
  s_{finite}=
  \begin{cases}
    1.0 & \text{si } r_{finite}\ge 0.95,\\
    0.5 & \text{si } 0.80\le r_{finite}<0.95,\\
    0.0 & \text{sinon.}
  \end{cases}
\]

Les spikes sont detectes sur les differences successives des valeurs finies :
\[
  d_j=y_{j+1}-y_j.
\]
Si $\sigma_d>10^{-12}$, le nombre de spikes est :
\[
  n_{spike}=\#\{j\mid |d_j|>6\sigma_d\}.
\]
La penalite est :
\[
  p_{spike}=\min\left(0.5,\frac{n_{spike}}{|\mathcal{W}|}\right).
\]
Le score canal est :
\[
  s_{canal}=\max(0,s_{finite}-p_{spike}).
\]
Le score global de qualite signal est la moyenne des scores canaux disponibles :
\[
  S_{signal}=\operatorname{mean}\{s_{canal,c}\}.
\]

\section{Residus physiques de l'index qualite}

Les residus suivants sont calcules sur toute l'acquisition, puis resumes par fenetre via moyenne, RMS, maximum absolu et quantile absolu a 95\%.

\subsection{Cinematique longitudinale}

\[
  r_{ax}(t)=a_x(t)-\frac{d v_x}{dt}(t).
\]

\subsection{Cinematique laterale}

Mode simple :
\[
  r_{ay,simple}(t)=a_y(t)-v_x(t)r(t).
\]

Mode avec vitesse laterale :
\[
  r_{ay}(t)=a_y(t)-\left(\frac{d v_y}{dt}(t)+v_x(t)r(t)\right).
\]

\subsection{Distance GPS}

Si une distance GPS curviligne $d_{gps}(t)$ est disponible, avec $v_{ref}=v$ si la vitesse scalaire est disponible, sinon $v_x$ :
\[
  r_{dist}(t)=\frac{d d_{gps}}{dt}(t)-v_{ref}(t).
\]

\subsection{Courbure}

Si le rayon GPS $R_{gps}$ est disponible :
\[
  \kappa_{gps}(t)=\frac{1}{R_{gps}(t)}.
\]
La courbure vehicule est :
\[
  \kappa_{veh}(t)=
  \begin{cases}
    \dfrac{r(t)}{\max(|v_x(t)|,5/3.6)} & \text{si } |v_x(t)|\ge 5/3.6,\\
    \mathrm{NaN} & \text{sinon.}
  \end{cases}
\]
Le residu est :
\[
  r_{\kappa}(t)=\kappa_{gps}(t)-\kappa_{veh}(t).
\]

\subsection{Roues}

Si les quatre vitesses de roues sont disponibles :
\[
  r_{wheel}(t)=\operatorname{std}\left(
  \omega_{fl}(t),\omega_{fr}(t),\omega_{rl}(t),\omega_{rr}(t)
  \right).
\]

\subsection{Calibration IMU}

\[
  r_{imu,x}(t)=a_{imu,x}(t)-a_x(t),
  \quad
  r_{imu,y}(t)=a_{imu,y}(t)-a_y(t),
  \quad
  r_{imu,z}(t)=a_{imu,z}(t)-a_z(t).
\]

\section{Residus GPS}

\subsection{Projection latitude/longitude en repere local}

Les latitudes et longitudes valides sont filtrees. Avec le premier point valide comme origine $(\varphi_0,\lambda_0)$, et $R_T=6\,371\,000~\mathrm{m}$ :
\[
  x_{gps}=R_T\cos(\varphi_0)(\lambda-\lambda_0),
  \qquad
  y_{gps}=R_T(\varphi-\varphi_0),
\]
ou les angles sont en radians.

\subsection{Vitesse GPS}

Les positions GPS sont lissees par moyenne mobile d'environ $0.5~\mathrm{s}$, puis derivees :
\[
  v_{gps,x}=\frac{d x_{gps}}{dt},\qquad
  v_{gps,y}=\frac{d y_{gps}}{dt},
\]
\[
  v_{gps}=\sqrt{v_{gps,x}^2+v_{gps,y}^2}.
\]
Les residus de vitesse GPS sont :
\[
  r_{gps,speed}=v_{gps}-v_{ref},
\]
\[
  r_{gps,body}=v_{gps}-\sqrt{v_x^2+v_y^2},
\]
ou, si $v_y$ est absent, $\sqrt{v_x^2+v_y^2}$ est remplace par $|v_x|$.

\subsection{Integration trajectoire vehicule}

L'angle de cap initial $\psi_0$ vient d'abord de la trajectoire GPS dans la fenetre si le deplacement total depasse $5~\mathrm{m}$, sinon du canal \texttt{gps.track}, sinon d'une derivee de trajectoire, sinon $0$.

Le yaw rate est integre :
\[
  \psi_i=\psi_0+\sum_{k=0}^{i} r_k\Delta t.
\]
La vitesse corps est projetee dans le repere monde :
\[
  \dot{x}_{veh}=v_x\cos\psi-v_y\sin\psi,
  \qquad
  \dot{y}_{veh}=v_x\sin\psi+v_y\cos\psi.
\]
Puis :
\[
  x_{veh,i}=x_{veh,0}+\sum_{k=1}^{i}\dot{x}_{veh,k}\Delta t,
  \qquad
  y_{veh,i}=y_{veh,0}+\sum_{k=1}^{i}\dot{y}_{veh,k}\Delta t.
\]

L'erreur de position brute est :
\[
  r_{gps,pos}(i)=
  \sqrt{(x_{gps,i}-x_{veh,i})^2+(y_{gps,i}-y_{veh,i})^2}.
\]

\subsection{Erreur GPS alignee}

Pour une fenetre locale, le code centre la trajectoire GPS et la trajectoire integree, puis cherche une rotation 2D de type Procrustes.
Si $p_i$ est la position GPS centree et $q_i$ la position integree centree :
\[
  d=\sum_i q_{x,i}p_{x,i}+q_{y,i}p_{y,i},
  \qquad
  c=\sum_i q_{x,i}p_{y,i}-q_{y,i}p_{x,i},
\]
\[
  \theta=\operatorname{atan2}(c,d).
\]
La trajectoire integree est tournee de $\theta$, recentree sur le GPS, puis l'erreur alignee vaut :
\[
  r_{gps,pos,aligned}(i)=\|p_i-\tilde{q}_i\|_2.
\]

\subsection{Yaw rate GPS}

Si \texttt{gps.track} existe, le cap GPS est unwrappe et derive :
\[
  r_{gps,yaw}(t)=\frac{d\psi_{gps}}{dt}(t)-r(t),
\]
uniquement lorsque la vitesse corps est superieure au seuil basse vitesse $5/3.6~\mathrm{m/s}$.

\section{Transformation residu vers score}

Chaque residu RMS peut etre transforme en score par :
\[
  s(e;\tau)=\frac{1}{1+\max(e,0)/\tau}.
\]
Un residu nul donne $1$. Un residu egal a $\tau$ donne $0.5$.

Les seuils $\tau$ utilises sont :

\begin{longtable}{ll}
\toprule
Score & Residus et seuils \\
\midrule
\texttt{kinematic\_score} &
$r_{dist}$ avec $\tau=1.5$, $r_{\kappa}$ avec $\tau=0.05$ \\
\texttt{acceleration\_score} &
$r_{ax}$ avec $\tau=1.0$, $r_{ay}$ et $r_{ay,simple}$ avec $\tau=1.5$ \\
\texttt{gps\_score} &
$r_{gps,pos}$ avec $\tau=10$, $r_{gps,pos,aligned}$ avec $\tau=5$,
$r_{gps,speed}$ et $r_{gps,body}$ avec $\tau=2$, $r_{gps,yaw}$ avec $\tau=0.2$ \\
\texttt{wheel\_score} &
$r_{wheel}$ avec $\tau=1.0$ \\
\texttt{calibration\_score} &
$r_{imu,x}$, $r_{imu,y}$, $r_{imu,z}$ avec $\tau=1.0$ \\
\bottomrule
\end{longtable}

Les moyennes ignorent les scores absents :
\[
  S_{phys}=\operatorname{mean}(S_{kin},S_{acc},S_{gps},S_{wheel},S_{calib}),
\]
\[
  S_{global}=\operatorname{mean}(S_{signal},S_{phys}).
\]

Important : dans le code actuel, $S_{global}$ est un score absolu par fenetre. Il n'est pas normalise par cluster. Le clustering ajoute un contexte de comparaison mais ne modifie pas ce score.

\section{Scores d'excitation et flags d'usage}

Les scores d'excitation sont :
\[
  S_{lat\_exc}=\min\left(1,\frac{\RMS(a_y)}{3}\right),
  \qquad
  S_{long\_exc}=\min\left(1,\frac{\RMS(a_x)}{2}\right).
\]

Les flags d'usage sont :
\[
  \texttt{usable\_for\_regression}
  =
  (S_{global}\ge 0.65)\land(p_{low}<0.5),
\]
\[
  \texttt{usable\_for\_lateral\_identification}
  =
  \texttt{usable\_for\_regression}\land(S_{lat\_exc}\ge 0.25),
\]
\[
  \texttt{usable\_for\_longitudinal\_identification}
  =
  \texttt{usable\_for\_regression}\land(S_{long\_exc}\ge 0.25),
\]
\[
  \texttt{usable\_for\_tire\_parameter\_estimation}
  =
  \texttt{usable\_for\_lateral\_identification}
  \land
  (r_{wheel}\text{ a le statut computed}),
\]
\[
  \texttt{usable\_for\_indicator\_computation}
  =
  (S_{signal}\ge 0.8).
\]

\section{Bins et familles de manoeuvre}

Les bins vitesse sont construits sur $\overline{v_x}^{km/h}$ :
\[
  [0,20),[20,40),[40,60),[60,80),[80,100),[100,120),
\]
sinon \texttt{out\_of\_range} ou \texttt{unknown}.

Les bins d'acceleration laterale :
\[
\begin{cases}
\texttt{faible} & \RMS(a_y)<1,\\
\texttt{moyen} & 1\le \RMS(a_y)<3,\\
\texttt{fort} & 3\le \RMS(a_y)<6,\\
\texttt{tres\_fort} & \RMS(a_y)\ge 6.
\end{cases}
\]

Les bins d'acceleration longitudinale :
\[
\begin{cases}
\texttt{faible} & \RMS(a_x)<0.5,\\
\texttt{moyen} & 0.5\le \RMS(a_x)<2,\\
\texttt{fort} & \RMS(a_x)\ge 2.
\end{cases}
\]

La direction de manoeuvre est determinee par le yaw rate :
\[
  f_+=\Pr(r>0.03),\qquad f_-=\Pr(r<-0.03).
\]
Si $f_+>0.2$ et $f_->0.2$, la direction est \texttt{mixed}. Sinon le signe de la moyenne de $r$ donne \texttt{left}, \texttt{right} ou \texttt{straight}.

La famille de manoeuvre suit les regles :
\[
\begin{array}{ll}
\texttt{stopped} & \overline{v_x}^{km/h}<1,\\
\texttt{straight} & \overline{v_x}^{km/h}\ge 5,\ \RMS(a_y)<0.5,\ \RMS(r)<0.03,\ \RMS(\delta)<0.02,\\
\texttt{lateral} & \RMS(a_y)\ge 1,\ \RMS(r)\ge 0.05,\ \RMS(a_x)<1,\\
\texttt{longitudinal} & \RMS(a_x)\ge 1,\ \RMS(a_y)<1,\\
\texttt{combined} & \RMS(a_x)\ge 1,\ \RMS(a_y)\ge 1,\\
\texttt{unknown} & \text{sinon.}
\end{array}
\]

L'identifiant de groupe de similarite est une concatenation :
\[
\texttt{family\_\_v\_speedBin\_\_ay\_latBin\_\_steer\_steerBin\_\_dir\_direction}.
\]

\section{Clustering de fenetres comparables}

Le script \texttt{build\_window\_cluster\_index.py} lit \texttt{window\_quality\_index.csv}.
Les features utilisees pour clusteriser sont :
\[
\begin{aligned}
  z = (&\overline{v_x}^{km/h},\RMS(a_x),\RMS(a_y),\RMS(r),\RMS(\delta),\\
       &p_{zero},p_{low},S_{lat\_exc},S_{long\_exc}).
\end{aligned}
\]
Seules les colonnes disponibles sont gardees. Il faut au moins trois features.

Le clustering est fait separement par valeur de \texttt{maneuver\_family}, si le groupe contient au moins \texttt{min\_group\_size}=30 fenetres.

\subsection{Standardisation robuste}

Pour chaque feature $x$ dans un groupe :
\[
  m=\median(x),\qquad
  q_{25}=Q_{0.25}(x),\qquad
  q_{75}=Q_{0.75}(x),
\]
\[
  s=q_{75}-q_{25}.
\]
Si $s$ est nul ou non fini, le code utilise l'ecart-type ; si celui-ci est aussi nul, $s=1$.
Les valeurs manquantes sont remplacees par la mediane, puis :
\[
  \tilde{x}_i=\frac{x_i-m}{s}.
\]

\subsection{DBSCAN}

Dans l'espace standardise $\tilde{z}$, le voisinage de rayon $\varepsilon=1.35$ est :
\[
  \mathcal{N}_{\varepsilon}(i)=\{j\mid \|\tilde{z}_j-\tilde{z}_i\|_2\le \varepsilon\}.
\]
Un point est coeur si :
\[
  |\mathcal{N}_{\varepsilon}(i)|\ge \texttt{min\_samples}=18.
\]
Les clusters sont les composantes connectees atteignables depuis les points coeurs. Les points non atteints restent \texttt{noise}.

\subsection{Distance au centre de cluster}

Pour chaque cluster $C$, le centre robuste est :
\[
  c_C=\median\{\tilde{z}_i\mid i\in C\}
\]
calcule composante par composante. La distance stockee est :
\[
  d_i=\|\tilde{z}_i-c_C\|_2.
\]

Le fichier \texttt{window\_cluster\_index.csv} reprend toutes les colonnes de \texttt{window\_quality\_index.csv} et ajoute notamment :
\[
  \texttt{cluster\_id},\quad
  \texttt{cluster\_label},\quad
  \texttt{cluster\_is\_noise},\quad
  \texttt{cluster\_size},\quad
  \texttt{cluster\_distance\_to\_center}.
\]

\section{Agregation par date et vitesse}

Le script \texttt{build\_speed\_day\_summary.py} groupe les fenetres par :
\[
  (\texttt{date},\texttt{speed\_bin}).
\]
Pour chaque colonne numerique $x$, il calcule :
\[
  \operatorname{mean}(x),\qquad
  \median(x),\qquad
  Q_{0.95}(x).
\]
Pour chaque flag booleen $b$, il calcule le ratio :
\[
  \operatorname{ratio}(b)=\frac{1}{n}\sum_{i=1}^{n}\mathbf{1}(b_i=\text{True}).
\]
Il stocke aussi le nombre de fenetres, le nombre de fichiers et la famille de manoeuvre dominante.

\section{Index de drift}

\texttt{build\_drift\_index.py} ne recalcule pas de nouvelle metrique physique.
Il lit les JSON exportes, extrait les features globales par canal et les regroupe dans \texttt{drift\_index.json}.
La logique mathematique est donc celle des features globales de la section precedente.

\section{Index d'anomalies de correlation}

Le script \texttt{build\_correlation\_anomaly\_index.py} exploite les moyennes de fenetres glissantes des JSON :
\[
  x_w=\operatorname{mean}_{\mathcal{W}}(x),\qquad
  y_w=\operatorname{mean}_{\mathcal{W}}(y).
\]

\subsection{Reference robuste par piste}

Les fichiers sont regroupes par \texttt{track\_id}, derive du nom de fichier.
Pour chaque paire de capteurs $(x,y)$ et chaque piste, le code collecte les couples $(x_w,y_w)$.

Les echelles robustes sont :
\[
  m_x=\median(x_w),\qquad
  s_x=1.4826\,\median(|x_w-m_x|),
\]
\[
  m_y=\median(y_w),\qquad
  s_y=1.4826\,\median(|y_w-m_y|).
\]
Les valeurs standardisees sont :
\[
  z_x=\frac{x_w-m_x}{s_x},\qquad
  z_y=\frac{y_w-m_y}{s_y}.
\]
La correlation de reference est :
\[
  \rho=\operatorname{corr}(z_x,z_y).
\]

Le residu relationnel est :
\[
  r_{rel}=\frac{z_y-\rho z_x}{\sqrt{\max(1-\rho^2,0.10)}}.
\]
Sa reference robuste est :
\[
  m_r=\median(r_{rel}),\qquad
  s_r=\max(1.4826\,\median(|r_{rel}-m_r|),0.50).
\]

\subsection{Score d'anomalie relationnelle}

Pour une fenetre donnee :
\[
  A=\frac{|r_{rel}-m_r|}{s_r}.
\]
Une anomalie candidate est gardee si :
\[
  A\ge 2.5.
\]
La severite est :
\[
\begin{cases}
\texttt{high} & A\ge 5.0,\\
\texttt{medium} & 3.5\le A<5.0,\\
\texttt{low} & A<3.5.
\end{cases}
\]

\section{Flux historique : export 01 et features 02}

\subsection{Export historique 01}

\texttt{01\_export\_dxd\_resampled\_json.py} selectionne seulement les fichiers qui possedent tous les canaux cibles, construit une grille a $100~\mathrm{Hz}$, interpole les canaux et exporte un JSON canonique. Il ne calcule pas les residus physiques avances de l'index qualite.

\subsection{Features univariees 02}

\texttt{02\_extract\_features.py} construit des fenetres de $1~\mathrm{s}$ avec un pas de $1~\mathrm{s}$ par defaut.
Les canaux sont gardes si leur ratio global de valeurs finies est au moins $0.90$ et s'il reste au moins trois canaux.

Dans chaque fenetre, les valeurs manquantes de chaque canal sont interpolees. Si un canal est entierement manquant dans une fenetre, la fenetre est rejetee.

Pour chaque canal $y$, les features sont :
\[
\begin{aligned}
&\overline{y},\ \sigma_y,\ \min y,\ \max y,\ \max y-\min y,\ \median(y),\\
&Q_{0.25}(y),\ Q_{0.75}(y),\ \IQR(y),\ \operatorname{mean}(|y|),\ \RMS(y).
\end{aligned}
\]
La pente lineaire est calculee par regression sur l'indice $i$ :
\[
  a=\frac{\sum_i (i-\overline{i})(y_i-\overline{y})}{\sum_i(i-\overline{i})^2}.
\]
Les differences donnent :
\[
  \sigma_{\Delta y}=\operatorname{std}(y_{i+1}-y_i),
  \qquad
  \max|\Delta y|=\max_i |y_{i+1}-y_i|.
\]
Le taux de changement de signe est :
\[
  ZCR=\frac{1}{n-1}\sum_{i=1}^{n-1}\mathbf{1}(\operatorname{sign}(y_{i+1})\ne \operatorname{sign}(y_i)).
\]
L'autocorrelation au lag 1 est :
\[
  \rho_1=
  \frac{\sum_i (y_i-\overline{y_0})(y_{i+1}-\overline{y_1})}
       {\sqrt{\sum_i(y_i-\overline{y_0})^2}\sqrt{\sum_i(y_{i+1}-\overline{y_1})^2}}.
\]

\section{Point important : qualite absolue versus qualite relative au cluster}

Le calcul actuel produit :
\[
  S_{global}=\operatorname{mean}(S_{signal},S_{phys})
\]
avant clustering. Ce score mesure une qualite absolue au sens des seuils de residus fixes.

Le clustering ne recalcule pas :
\[
  S_{global\mid cluster}
\]
et ne produit pas de percentile ou de z-score robuste des residus dans le cluster.

Si l'objectif est de comparer une fenetre uniquement a des fenetres de dynamique similaire, une extension mathematiquement naturelle serait, pour chaque cluster $C$ et chaque residu $e$ :
\[
  z_{e,i}^{(C)}=
  \frac{e_i-\median(e_j:j\in C)}
       {\max(1.4826\,\median(|e_j-\median_C(e)|:j\in C),\epsilon)}.
\]
On pourrait ensuite definir un score relatif :
\[
  S_{cluster,i}
  =
  \frac{1}{1+\operatorname{mean}_{e\in\mathcal{E}} \max(z_{e,i}^{(C)},0)}
\]
ou bien un percentile de qualite dans le cluster. Cette extension n'est pas presente dans le code actuel.

\end{document}
